\subsection{Thiết kế và hiện thực phân hệ lõi luật chơi}

\subsubsection{Tổ chức kiến trúc lõi tách rời và cấu hình cân bằng tập trung}
Từ phân tích đánh đổi giữa các hướng triển khai trên Unity ở Mục 3.2, toàn bộ phần xử lý luật chơi của \textit{Convergence Chess} được thiết kế tách rời hoàn toàn khỏi đồ họa 3D và tổ chức quanh ba khối thành phần chính (minh họa tổ chức đối tượng và tài nguyên cấu hình trong Unity tại Hình \ref{fig:unity_core_config}):
\begin{itemize}
    \item \textbf{Nguồn cấu hình cân bằng tập trung:} Lưu toàn bộ tham số trò chơi (ngân sách $12$ điểm đội hình, chi phí điểm và chi phí AP của từng loại quân, mức hồi $+2$ AP mỗi lượt, trần tích lũy $5$ AP, phụ phí hợp thể hiện đặt bằng $0$, giới hạn $3$ lượt bỏ và $60$ lượt đấu) trong một tài nguyên dữ liệu tĩnh độc lập (\texttt{ScriptableObject} của Unity). Cách tổ chức này giúp nhóm điều chỉnh cân bằng lối chơi nhanh chóng mà không phải sửa trực tiếp logic xử lý.
    \item \textbf{Cấu trúc trạng thái ván đấu thống nhất:} Lưu tập trung toàn bộ dữ liệu của một trận đấu trong cùng một cấu trúc duy nhất, gồm vị trí quân cờ trên lưới $36$ ô, $1$--$2$ vật cản vùng trung tuyến (ô $12$ đến ô $23$) sinh theo hạt giống xác định, tập $9$ kỹ năng nội tại của từng quân cờ (được biểu diễn gọn dưới dạng tập cờ trạng thái giúp thao tác hợp thể và sao chép bàn cờ diễn ra tức thời), điểm hành động hai bên, trạng thái đòn bồi Surge và số lượt còn lại.
    \item \textbf{Cô lập hoàn toàn theo từng bàn đấu:} Thay vì dùng bộ quản lý toàn cục dễ gây xung đột dữ liệu, mỗi bàn đấu sở hữu một vùng quản lý dịch vụ riêng được lắp ráp linh hoạt theo từng chế độ chơi (đấu với máy, đấu hai người cùng máy, đấu trực tuyến hoặc mô phỏng huấn luyện AI). Ở chế độ huấn luyện AI, tầng diễn hoạt đồ họa được tắt hoàn toàn để chạy song song hàng chục bàn cờ cùng lúc.
\end{itemize}

\begin{figure}[H]
    \centering
    \screenshotplaceholder{Content/Core và AI/images/unity\_core\_config\_hierarchy.png}{Chụp ảnh cấu trúc bàn đấu mẫu (Prefab Arena) trong cửa sổ Hierarchy cùng tài nguyên cấu hình cân bằng trò chơi (ScriptableObject) trong Unity Editor.}
    \caption{Tổ chức bàn đấu độc lập và nguồn cấu hình cân bằng tập trung trong Unity 6}
    \label{fig:unity_core_config}
\end{figure}

\subsubsection{Luồng vận hành một lượt đấu và kiểm định luật chơi}
Vòng lặp trận đấu vận hành theo cơ chế bất đồng bộ từng lượt, kết nối với bên ngoài qua hai cổng thống nhất (cổng tiếp nhận hành động và cổng phát sự kiện). Một lượt chơi diễn ra tuần tự qua bốn bước:
\begin{enumerate}
    \item \textbf{Khởi tạo đầu lượt:} Hệ thống cộng điểm hành động đầu lượt ($+2$ AP cơ bản và cộng thêm từ kỹ năng đặc biệt, tối đa $5$ AP) và kiểm tra xem phe tới lượt còn ít nhất một nước đi hoặc quyền bỏ lượt hợp lệ hay không; nếu không còn hành động nào, trận đấu kết thúc và xử thua cho phe đó.
    \item \textbf{Tiếp nhận nước đi theo chế độ chơi:} Hệ thống chờ nhận lệnh hành động (di chuyển, tấn công, hợp thể hoặc bỏ lượt) từ đúng nguồn điều khiển của phe đang tới lượt, bao gồm thao tác chuột của người chơi, dữ liệu từ đối thủ trực tuyến hoặc quyết định từ tác nhân AI.
    \item \textbf{Thẩm định toàn diện tính hợp lệ:} Bộ thẩm định luật chơi kiểm tra đồng thời mọi điều kiện của nước đi:
    \begin{itemize}
        \item Đủ lượng AP cần thiết theo loại quân và khoảng cách di chuyển; đường trượt không bị chặn bởi quân cờ khác hoặc vật cản trung tuyến (ngoại trừ quân có kỹ năng nhảy chữ L).
        \item Tuân thủ đầy đủ quy tắc hợp thể: cấm hợp thể hai quân cùng loại gốc, mỗi quân chỉ tham gia hợp thể tối đa một lần trong trận và không cho phép quân Vua tham gia hợp thể.
        \item Khi đang ở giai đoạn đòn bồi Surge, chỉ cho phép quân cờ vừa hợp thể thực hiện hành động hoặc chọn kết thúc lượt để giữ lại AP.
        \item Kiểm tra trạng thái chiếu Vua, xử lý tương tác giữa Khiên hộ mệnh với Phá khiên và phong cấp Tốt khi tiến tới hàng cuối sân đối phương.
    \end{itemize}
    \item \textbf{Cập nhật trạng thái và phát đợt sự kiện:} Sau khi nước đi hợp lệ được áp dụng lên trạng thái bàn cờ, lõi xử lý gom toàn bộ các hệ quả phát sinh trong nước đi đó (di chuyển, hợp thể, vỡ khiên, tiêu diệt quân, phong cấp Tốt, kích hoạt Surge, tiêu hao AP, đổi lượt hoặc kết thúc trận) thành một đợt sự kiện tuần tự gửi sang tầng hiển thị. Riêng trong quá trình duyệt cây tìm kiếm của AI, lõi luật hỗ trợ đi thử và hoàn tác trực tiếp mà không tạo sự kiện hình ảnh để tối đa hóa tốc độ mô phỏng.
\end{enumerate}

\subsubsection{Đánh giá ưu và nhược điểm của kiến trúc lõi tách rời}
\begin{itemize}
    \item \textbf{Ưu điểm:} Bảo đảm tính nhất quán tuyệt đối giữa mọi chế độ chơi (đấu máy, đấu hai người cùng máy và đấu trực tuyến); đạt tốc độ mô phỏng hàng trăm ván đấu trong vài chục giây cho huấn luyện AI nhờ chạy không đồ họa và biểu diễn trạng thái gọn nhẹ; thuận lợi khi viết kiểm thử tự động cho các tình huống biên của luật chơi.
    \item \textbf{Nhược điểm và đánh đổi:} Chi phí thiết kế ban đầu cao hơn so với cách viết trực tiếp vào đối tượng 3D của Unity; mỗi khi bổ sung một cơ chế luật mới, nhóm phải đồng bộ thêm loại sự kiện tương ứng ở lớp cầu nối sang tầng hiển thị.
\end{itemize}

\subsection{Thiết kế và hiện thực hệ thống trí tuệ nhân tạo}

Đóng góp nghiên cứu trọng tâm của Giai đoạn 1 là xây dựng và kiểm chứng hai dòng tác nhân AI chính là \textbf{Hybrid Bot} và \textbf{LLM Bot}, cùng giao tiếp với lõi trò chơi thông qua một giao diện nhận nước đi thống nhất. Để đánh giá khách quan, báo cáo đối chiếu hai dòng này với hai mốc truyền thống là mạng tự quyết định nước đi (không tìm kiếm) và Minimax thuần ở Chương 6.

\subsubsection{Bộ mã hóa trạng thái bàn cờ}
Toàn bộ hệ thống AI dùng chung một bộ mã hóa trạng thái gồm \textbf{654 đặc trưng số thực} ($36\text{ ô} \times 18\text{ kênh} + 6\text{ biến toàn cục}$), luôn được chuẩn hóa theo góc nhìn của phe đang tới lượt:
\begin{itemize}
    \item Mỗi ô trên bàn cờ $6 \times 6$ được mô tả bằng $18$ kênh thông tin: $1$ kênh đánh dấu ô vật cản, $6$ kênh loại quân cờ có dấu ($+1$ với quân mình, $-1$ với quân địch cho 6 loại quân từ Tốt đến Vua), $9$ kênh tương ứng với $9$ kỹ năng nội tại có dấu theo phe sở hữu, $1$ kênh trạng thái đã từng hợp thể và $1$ kênh đánh dấu ô đang kích hoạt Surge.
    \item Sáu biến toàn cục biểu diễn lượng AP hiện có của hai bên, số lượt bỏ lượt còn lại của hai bên, tiến độ trận đấu so với giới hạn $60$ lượt và cờ báo hiệu đang ở giai đoạn Surge.
\end{itemize}
Không gian hành động rời rạc gồm $36 \times 36 + 1 = 1{,}297$ lựa chọn (ánh xạ mọi cặp ô xuất phát -- ô đích và hành động cuối cùng là bỏ lượt). Toàn bộ không gian hành động được lọc chặt qua bộ thẩm định luật chơi ở từng lượt để tác nhân AI không bao giờ chọn nước đi phạm luật.

\subsubsection{Hướng tác nhân 1: Hybrid Bot}
\textbf{Hybrid Bot} kết hợp giữa tìm kiếm cây trò chơi cổ điển và học sâu nhằm khắc phục cả tầm nhìn ngắn hạn của mạng chính sách thuần túy lẫn giới hạn của hàm đánh giá viết tay. Hybrid Bot gồm ba thành phần:
\begin{enumerate}
    \item \textbf{Khung tìm kiếm Minimax Alpha-Beta tối ưu:} Cài đặt tìm kiếm cây trò chơi với cắt tỉa Alpha-Beta, bảng hoán vị 64-bit Zobrist, ghi nhớ nước cắt tỉa tốt theo độ sâu, bảng điểm lịch sử cho nước đi tĩnh, ưu tiên duyệt nước ăn quân hoặc hợp thể và tự động gia hạn tối đa $2$ tầng tìm kiếm khi xuất hiện đòn bồi Surge. Cấu hình mặc định tìm kiếm ở độ sâu $4$.
    \item \textbf{Mạng chính sách gợi ý nước đi:} Huấn luyện bằng học tăng cường PPO tự đấu \cite{ppo} trên nhiều bàn cờ chạy song song, kết hợp định hình phần thưởng theo hàm thế năng của mạng NNUE \cite{potential_shaping}:
    \[
        \Delta V_{\text{NNUE}} = \text{hệ số} \cdot \left(\gamma \Phi(s_{t+1}, p) - \Phi(s_t, p)\right)
    \]
    Trong trận đấu, mạng chính sách cung cấp phân phối xác suất ưu tiên nước đi: bộ tìm kiếm duyệt đủ độ sâu với nhóm khoảng $5$ nước đi được ưu tiên (hoặc nhóm đạt ngưỡng xác suất tích lũy $90\%$) và toàn bộ nước ăn quân hoặc hợp thể, còn các nước tĩnh khác chỉ được thăm dò ở độ sâu giảm $2$ tầng và được kiểm chứng lại ở độ sâu đầy đủ nếu xuất hiện điểm vượt trội.
    \item \textbf{Mạng đánh giá nhanh dạng dư (Residual NNUE):} Chấm điểm nhanh các thế cờ ở nút lá bằng mạng ba tầng ($654 \to 128 \to 64 \to 1$) \cite{nnue} khai thác tính thưa của bàn cờ $6 \times 6$. Điểm đánh giá viết tay đóng vai trò mốc neo, mạng nơ-ron chỉ học phần hiệu chỉnh bị giới hạn trong khoảng $[-3.0, +3.0]$. Giữa mạng chính sách (sinh dữ liệu tự đấu qua nhiều thế hệ) và mạng đánh giá (định hình phần thưởng ngược lại cho mạng chính sách) hình thành vòng lặp đồng huấn luyện nâng dần sức mạnh qua từng thế hệ.
\end{enumerate}
Kết quả thực nghiệm (Chương 6) cho thấy phần đóng góp rõ rệt nhất đến từ Residual NNUE kết hợp tìm kiếm; riêng phần mạng chính sách trong Hybrid Bot hiện chưa mang lại lợi ích vượt trội so với cấu hình chỉ gồm Minimax và Residual NNUE ở cùng độ sâu.

\subsubsection{Hướng tác nhân 2: LLM Bot (tiến hóa tập luật biểu tượng)}
Song song với Hybrid Bot dựa trên mạng nơ-ron, nhóm phát triển \textbf{LLM Bot} theo hướng trí tuệ nhân tạo biểu tượng có thể giải thích được (lấy cảm hứng từ FunSearch \cite{funsearch}):
\begin{itemize}
    \item Toàn bộ trạng thái bàn cờ $6 \times 6$ được biểu diễn gọn bằng các mặt nạ bit $36$-bit siêu nhẹ, cho phép Gemini API sinh và tiến hóa từng mô-đun quy tắc đánh giá ngắn gọn ($15$--$35$ dòng mã) ứng với $8$ ngữ cảnh chiến thuật riêng biệt: luôn áp dụng, khai cuộc chưa hợp thể, trung cuộc đã hợp thể, tàn cuộc ít quân, tranh chấp khiên hộ mệnh, Vua bị đe dọa, đua phong cấp Tốt và trạng thái tích trữ AP cao. Mỗi mô-đun đều được bọc cơ chế đánh giá đối xứng giữa hai phe để bảo đảm công bằng.
    \item Mỗi quy tắc mới do LLM đề xuất đều phải vượt qua \textbf{4 cổng kiểm định tự động}: (1) kiểm tra cú pháp an toàn và không gây cấp phát bộ nhớ động; (2) biên dịch thực tế kèm tối đa $1$ lượt tự sửa lỗi nếu biên dịch thất bại; (3) đo tốc độ thực thi và kiểm chứng tính đối xứng phản chiếu giữa hai phe; (4) thi đấu đối chứng cặp (đảo vai đi trước/đi sau trên cùng hạt giống bàn cờ) kết hợp lọc loại trừ để giữ lại các mô-đun thực sự mang lại lợi ích. Sau cắt tỉa, ngân hàng luật hiện tại giữ lại $4$ mô-đun: tận dụng đòn bồi Surge cùng AP dư để tạo mối đe dọa lên Vua và quân giá trị cao, tiến Tốt kiểm soát trung tâm, đánh giá an toàn vùng quanh Vua, và săn quân giá trị cao đang không được bảo vệ.
\end{itemize}

\subsubsection{Đánh giá ưu và nhược điểm của Hybrid Bot và LLM Bot}
Bảng \ref{tab:ai_pros_cons} tổng hợp trực diện ưu điểm và hạn chế của hai dòng tác nhân AI chính do nhóm triển khai:

\begin{table}[H]
    \centering
    \caption{So sánh ưu điểm và nhược điểm của Hybrid Bot và LLM Bot}
    \label{tab:ai_pros_cons}
    \begin{tabularx}{\textwidth}{L{2.4cm} Y Y}
        \toprule
        \textbf{Dòng tác nhân} & \textbf{Ưu điểm nổi bật} & \textbf{Nhược điểm và đánh đổi} \\
        \midrule
        \textbf{Hybrid Bot} & Sức mạnh chiến thuật cao nhờ kết hợp tìm kiếm và Residual NNUE (cấu hình Minimax cùng Residual NNUE thắng $51\%$ trước Minimax độ sâu 4 với khoảng $66$ mili-giây mỗi nước); có vòng lặp đồng huấn luyện đa thế hệ giúp tự cải thiện. & Quy trình huấn luyện nhiều tầng phức tạp; quyết định của mạng nơ-ron mang tính ``hộp đen'' khó giải thích bằng lời; thế hệ 8 tốn khoảng $200$ mili-giây mỗi nước ở độ sâu 4; ở thực nghiệm hiện tại phần mạng chính sách chưa vượt cấu hình Minimax cùng Residual NNUE thuần; cần huấn luyện lại khi đổi mạnh tham số cân bằng. \\
        \textbf{LLM Bot} & Tri thức đánh giá ở dạng tập luật biểu tượng minh bạch, đọc hiểu và kiểm chứng được; không tốn chi phí gọi API khi thi đấu vì luật đã biên dịch sẵn; rất nhanh (khoảng $13$ mili-giây mỗi nước) mà vẫn thắng $42\%$ trước Minimax độ sâu 4; dễ bật/tắt hoặc tinh chỉnh theo từng ngữ cảnh chiến thuật. & Quá trình tiến hóa luật phụ thuộc chất lượng gợi ý và chi phí gọi API ngoại tuyến; sức mạnh trần khó theo kịp mạng đánh giá học sâu trên các thế cờ phức tạp; mới có $4$ mô-đun được giữ lại; chưa có đối đầu trực tiếp với Hybrid Bot nên chưa thể kết luận hơn kém giữa hai hướng. \\
        \bottomrule
    \end{tabularx}
\end{table}

\subsection{Thiết kế và hiện thực phân hệ tùy biến nhân vật 3D (Smart Customization)}

Phân hệ tùy biến nhân vật 3D (Smart Customization) được thiết kế theo định hướng \textbf{tự chủ thuật toán khi chạy và sẵn sàng thương mại hóa}: ở bản phát hành, quy trình chạy cục bộ trên máy trạm (NVIDIA GeForce RTX 4050 Laptop, bộ nhớ đồ họa 6 GB) và áp trực tiếp lên bộ khung nhân vật chuẩn giải phẫu \textbf{MPFB v2} \cite{mpfb2} (Blender 4.2) trong Unity 6, hoàn toàn không gọi HRN hay IC-Light.

\subsubsection{Nhánh tái tạo hình học khuôn mặt 3D}
Thế hệ đầu của nhóm dự đoán $25$ thông số hình thái từ dữ liệu ảnh phối cảnh tổng hợp; thế hệ hiện tại đã thay thế hoàn toàn bằng \textbf{mạng học sinh hai nhánh} dự đoán $60$ thông số hình thái 3D, huấn luyện bằng chưng cất tri thức từ ảnh chân dung thật:
\begin{enumerate}
    \item \textbf{Giáo viên ngoại tuyến:} Mô hình HRN \cite{hrn} dựng hình học khuôn mặt 3D không biểu cảm từ ảnh chân dung; một bộ giải bình phương tối thiểu có ràng buộc chuyển hình học đó thành $60$ thông số hình thái của MPFB v2 (khoảng $0.3$ mili-giây mỗi khuôn mặt). Đây là khâu duy nhất còn dùng HRN và chỉ chạy ngoại tuyến để sinh nhãn.
    \item \textbf{Dữ liệu huấn luyện:} $5{,}000$ ảnh chân dung thật từ tập FairFace \cite{fairface} cân bằng $7$ nhóm sắc tộc ($2{,}500$ mẫu nam và $2{,}500$ mẫu nữ), mỗi ảnh được căn chỉnh góc xoay và cắt vuông $224 \times 224$.
    \item \textbf{Mạng học sinh hai nhánh:} Nhánh thị giác dùng ResNet34 \cite{resnet} đọc ảnh khuôn mặt để nắm bắt độ sâu, độ nhô của mũi, gò má, cằm; nhánh hình học dùng mạng nhiều lớp đọc $968$ đặc trưng từ điểm mốc MediaPipe \cite{mediapipe} (gồm cả các tỷ lệ phối cảnh). Hai nhánh được hòa trộn qua tầng có cổng, và một đầu phụ ước lượng tư thế đầu chỉ dùng khi huấn luyện. Mạng được huấn luyện $40$ chu kỳ với hàm mất mát kết hợp sai số đỉnh 3D tính bằng mili-mét, sai số Smooth L1 và thành phần bảo toàn hướng cùng biên độ để tránh dự đoán co về $0$, rồi xuất sang định dạng ONNX.
    \item \textbf{Giới hạn an toàn giải phẫu:} Mỗi thông số bị kẹp trong biên độ riêng; riêng vùng môi được khóa ở dáng khép tự nhiên để tránh hở khe môi lộ răng.
\end{enumerate}
$60$ thông số bao gồm: $9$ thông số sọ và trán, $3$ thông số cung mày, $11$ thông số mắt và mí, $20$ thông số mũi, $4$ thông số gò má, $5$ thông số môi ở dáng khép và $8$ thông số hàm và cằm.

\subsubsection{Nhánh xử lý bề mặt da, vật liệu mắt và hệ thống tóc 3D}
\begin{itemize}
    \item \textbf{Chiếu ngược kết cấu da và khử bóng đối xứng tự cài đặt:} Tham khảo ý tưởng chiếu tam giác từ nghiên cứu \textbf{MeInGame} \cite{meingame}, nhóm tự hiện thực phép chiếu ngược trên \textbf{PyTorch3D} \cite{pytorch3d} để chuyển màu da từ ảnh chân dung sang bản đồ bề mặt của nhân vật 3D, sau đó áp dụng thuật toán khử bóng và khử che khuất theo tính đối xứng hai bên khuôn mặt trong không gian màu Lab. Cách làm này trung hòa hiệu quả bóng đổ gắt hoặc tóc mái che một bên trong khi vẫn giữ được lỗ chân lông, lông mày và râu, không phụ thuộc mô hình khử bóng khuếch tán IC-Light \cite{iclight}.
    \item \textbf{Tách hốc mắt, khử vệt răng và hòa trộn biên cổ:} Hệ thống dùng mặt nạ tách bỏ vùng hốc mắt để giữ nguyên cấu trúc nhãn cầu 3D có khả năng xoay độc lập; áp dụng bộ lọc nhận diện men răng để loại bỏ vệt răng trắng nếu ảnh gốc đang cười; cân bằng sắc tố da trong không gian màu Lab kết hợp hòa trộn tháp đa tần số và cửa sổ chuyển tiếp mượt ($64$ điểm ảnh) giúp xóa viền ghép giữa mặt và cổ.
    \item \textbf{Vật liệu mắt phản sáng và hệ thống 14 kiểu tóc 3D:} Bề mặt mống mắt và củng mạc được cấu hình độ bóng cao ($0.96$) kèm lớp phủ màng nước mắt tạo đốm sáng tự nhiên trên giác mạc, đồng thời chuẩn hóa ngưỡng sáng tối thiểu của mống mắt ở mức $0.36$ để vân mắt luôn rõ nét. Hệ thống tóc gồm $14$ kiểu tóc 3D dùng chung cho cả hai giới cùng chế độ không tóc ($15$ lựa chọn), được gắn vào đầu nhân vật theo từng giới tính và tự động lấy mẫu màu tóc từ vùng viền trán của ảnh chân dung đầu vào.
\end{itemize}

\subsubsection{Đánh giá ưu và nhược điểm của giải pháp Smart Customization tự xây dựng}
\begin{itemize}
    \item \textbf{Ưu điểm:} Khi chạy thực tế không phụ thuộc mô hình phi thương mại nào; bảo toàn toàn bộ hệ khung xương hoạt họa của MPFB v2; giữ trung thực nét nhận dạng khuôn mặt gốc; tốc độ xử lý khoảng $5$ giây mỗi ảnh; so với thế hệ đầu, $60$ thông số 3D bao phủ thêm cấu trúc mũi, gò má, hàm và cằm nên khuôn mặt tái tạo có chiều sâu hơn.
    \item \textbf{Nhược điểm:} Nhãn huấn luyện vẫn xuất phát từ giáo viên HRN phi thương mại nên còn rủi ro pháp lý ở khâu dữ liệu, cần thay thế ở Giai đoạn 2; sai số hình học ghi nhận được đo so với giáo viên chứ chưa đo so với số liệu thật độc lập; do không dùng bản đồ chi tiết nếp nhăn từ HRN khi chạy nên vi chi tiết trên da ít nổi bật; chất lượng còn nhạy với ảnh chụp nghiêng góc lớn hoặc ngược sáng cực đoan; vẫn cần cạc đồ họa rời hoặc dịch vụ máy chủ.
\end{itemize}
